\documentclass[12pt,a4paper,openany,oneside]{book}

% --- Paquetes ---
\usepackage[utf8]{inputenc}
\usepackage[spanish, es-tabla]{babel}
\usepackage{amsmath, amsfonts, amssymb}
\usepackage{graphicx}
\usepackage{geometry}
\usepackage{hyperref}
\usepackage{booktabs}
\usepackage{fancyhdr}
\usepackage{listings}
\usepackage{xcolor}
\usepackage{tikz}
\usepackage{enumitem}
\usepackage{array}
\usetikzlibrary{arrows.meta, positioning, shapes.geometric}

\geometry{margin=2.5cm}

% --- Estilo de Código Python ---
\definecolor{codegreen}{rgb}{0,0.6,0}
\definecolor{codegray}{rgb}{0.5,0.5,0.5}
\definecolor{codepurple}{rgb}{0.58,0,0.82}
\definecolor{backcolour}{rgb}{0.95,0.95,0.92}

\lstset{
    language=Python,
    backgroundcolor=\color{backcolour},   
    commentstyle=\color{codegreen},
    keywordstyle=\color{blue},
    numberstyle=\tiny\color{codegray},
    stringstyle=\color{codepurple},
    basicstyle=\ttfamily\small,
    breaklines=true,
    frame=single,
    showstringspaces=false
}

% --- Cabecera y Pie de Página ---
\pagestyle{fancy}
\fancyhf{}
\renewcommand{\headrulewidth}{0.4pt}
\renewcommand{\footrulewidth}{0.4pt}
\fancyhead[L]{\small Carlos César Sánchez Coronel}
\fancyhead[R]{\small Machine Learning para IA}
\fancyfoot[L]{\small Carlos César Sánchez Coronel}
\fancyfoot[C]{\thepage}

% --- Portada ---
\title{
    \textbf{Sílabo del Curso} \\ 
    \Huge \textbf{MACHINE LEARNING} \\
    \Large \textsc{para Inteligencia Artificial} \\[1cm]
    \normalsize Fundamentos, algoritmos y aplicaciones prácticas
}
\author{
    \small Por \\ \vspace{0.2cm}
    \Large \textsc{Carlos César Sánchez Coronel}
}
\date{2026}

\begin{document}

\maketitle
\tableofcontents

% ------------------------------------------------------------------
% PRESENTACIÓN DEL CURSO
% ------------------------------------------------------------------
\chapter*{Presentación del Curso}
\addcontentsline{toc}{chapter}{Presentación del Curso}

Este curso está diseñado para futuros ingenieros e científicos de datos que deseen dominar los fundamentos del machine learning desde una perspectiva práctica y orientada a la industria. A lo largo de 14 semanas, exploraremos los algoritmos más importantes, técnicas de preprocesamiento, evaluación, interpretabilidad y despliegue, con un enfoque en las preguntas y habilidades que se esperan en entrevistas técnicas y en el trabajo diario.

Cada Modelo de ML contiene:
\begin{itemize}
    \item \textbf{Problematica de negocio}: Problematica de negocio.
    \item \textbf{Modelado de ML}: Selección de modelo
    \item \textbf{Fundamento matematico y computacional}: Matematica y computación que incluye el modelo.
    \item \textbf{Pipeline de datos}: Obtención de datos, preprocesamiento, feature engineering, evaluación.
    \item \textbf{Análisis de metricas}: Interpretacion de metricas del modelo.
    \item \textbf{Comunicación}: Storytelling, elevator pitch a equipo tecnico y no tecnico.


% ------------------------------------------------------------------
% SEMANA 1: INTRODUCCIÓN
% ------------------------------------------------------------------
\chapter{Semana 1: Introducción al ML}

\section{Logro de la sesión}
Comprender el flujo completo de un proyecto de machine learning, identificar los tipos de aprendizaje y reconocer casos de uso en la industria.

\subsection{Conceptos}
\begin{itemize}
    \item Definición de Machine Learning, Deep Learning y su relación con la IA.
    \item Problemas de negocios y Tipos de aprendizaje: supervisado, no supervisado, por refuerzo.
    \item Metodos parametricos y no parametricos en ML.
    \item Historia y auge de los Modelos de ML.
    \item Ciclo de vida y metodologia CRISP-DM.
    \item Diferencias: Data engineer, Data Scientist, ML Engineer, MLOps.
\end{itemize}

\subsection{Aplicaciones y casos típicos}
\begin{itemize}
    \item Sistemas de recomendación (Netflix, Amazon).
    \item Detección de fraude en transacciones bancarias.
    \item Mantenimiento predictivo en maquinaria industrial.
    \item Clasificación de imágenes (diagnóstico médico, vehículos autónomos).
    \item Predicción de series temporales (ventas, demanda energética).
\end{itemize}

\subsection{Fundamentos matematicos y computacionales}
\begin{itemize}
    \item Matemáticas detrás de los ML clasico: Funciones de costos, optimización, generalizado y especifico por modelos de regresion, clasificacion y clustering.
    \item Complejidad algoritmica que incluyen los modelos.
    \item Hardware: CPU, GPU, TPU.
\end{itemize}

\subsection{Métricas}
\begin{itemize}
    \item Tipos de metricas en ML e importancia de la métrica en el negocio.
\end{itemize}

\subsection{Laboratorio: Ver Colab}
\begin{itemize}
    \item Repaso de Numpy
    \item Repaso de Pandas
    \item Repaso de Matplotlib
    \item Repaso de Seaborn
\end{itemize}

\subsection{Reto: 1 punto}
\begin{itemize}
    \item Escoger una empresa tech famosa, y detectar la mayor cantidad de modelos de ML que puede tener. Ejm. en Uber, Netflix, Spotify, TikTok, etc.
\end{itemize}

% ------------------------------------------------------------------
% SEMANA 2: ANÁLISIS EXPLORATORIO Y FEATURE ENGINEERING
% ------------------------------------------------------------------
\chapter{Semana 2: Análisis Exploratorio y Feature Engineering}

\section{Logro de la sesión}
Realizar un análisis exploratorio y aplicar técnicas de feature engineering para preparar datos de calidad.

\subsection{Conceptos}
\begin{itemize}
    \item Importancia del EDA
    \item Casos famosos de exito y fracaso en negocios
    \item Pipeline de EDA
    \item Limpieza de datos (valores faltantes, valores duplicados, errores de tipo de dato, tratamiento basicode outliers, inconsistencias y errores logicos)
    \item Analisis univariado (Variables numericas: medidas de tendencia central, medidas de dispersion, medidas de forma, tipos de visualizaciones; Variables categoricas: frecuencias, moda, proporciones, tipos de visualizaciones; Variables de Fecha: Componentes de fechas, tipos de visualizaciones)
    \item Analisis bivariado (Variables numericas vs numericas, categoricas vs categoricas, numericas vs categoricas)
    \item Analisis multivariado (Matrices de correlacion, Pairplots, Heatmaps)

    \item Importancia del Feature engineering
    \item Pipeline de feature engineering
    \item Creación de nuevas características (a partir de numericas, categoricas, fechas, agregaciones, texto,etc)
    \item Codificación de variables categoricas (One-Hot Encoding, Label Encoding, Target Encoding, Binary Encoding, Hashing Encoding)
    \item Escalado y Normalizacion (Z-score, Min-max, RobustScaler, Normalizer)
    \item Manejo de datos desbalanceados (Oversampling, Undersampling, SMOTE)
    \item Imputación de datos (Mean, Median, Mode, KNNImputer, IterativeImputer)
    \item Manejo de outliers (variable indicadora de outlier, winsorizacion, transformaciones para reducir impacto)
    \item Selección de características (Filtros, Wrappers, Embebidos)
    \item Variables derivadas del negocio (ratios, agregaciones temporales, interacciones)
    
\end{itemize}

\subsection{Comunicacion}
\begin{itemize}
    \item Storytelling e insights con data dummy
    \item Elevator pitch a equipo tecnico y no tecnico
\end{itemize}

\subsection{Laboratorio: Ver Colab}
\begin{itemize}
    \item EDA
    \item Introducción a scikit-learn
    \item Feature engineering
\end{itemize}

\subsection{Reto: 1 punto}
\begin{itemize}
    \item Buscar 3 casos de exito del EDA y/o feature engineering en empresas big tech.
\end{itemize}

\subsection{Anexo: Fundamentos matematicos y computacionales}
\begin{itemize}
    \item Estadistica descriptiva
    \item Probabilidad
    \item Estadistica inferencial
    \item Optimizacion computacional
\end{itemize}

% ------------------------------------------------------------------
% SEMANA 3: REGRESIÓN LINEAL Y REGULARIZACIÓN
% ------------------------------------------------------------------
\chapter{Semana 3: Regresión Lineal y Regularización}

\section{Logro de la sesión}
Construir e interpretar modelos de regresión lineal, aplicando regularización y sus metricas.

\subsection{Problematica de negocio}
\begin{itemize}
    \item Tipos y ejemplos de problemas de regresion
    \item Solucion: regresión simple y multiple
\end{itemize}

\subsection{Modelado}
\begin{itemize}
    \item Requisitos del modelo
    \item Regresión lineal simple y múltiple: formulación, interpretación de coeficientes.
    \item Supuestos del modelo: linealidad, independencia, homocedasticidad, normalidad de residuos.
    \item Regularización: Ridge (L2), Lasso (L1) y Elastic Net.
    \item Interpretación de la regularización y selección de hiperparámetros.
    \item Plantilla base de Python: Imports, Split de datos, Modelo, Entrenamiento, Predicciones, Evaluacion
\end{itemize}

\subsection{Métricas}
\begin{itemize}
    \item Error Cuadrático Medio (MSE), Raíz del Error Cuadrático Medio (RMSE), Error Absoluto Medio (MAE), Coeficiente de determinación \(R^2\).
    \item Interpretacion y selección de metricas
\end{itemize}

\subsection{Comunicacion}
\begin{itemize}
    \item Storytelling e insights con data dummy
    \item Elevator pitch a equipo tecnico y no tecnico
\end{itemize}

\subsection{Reto: 1 punto}
\begin{itemize}
    \item Investigar el computo necesario para entrenar un modelo de regresión lineal en un dataset de 10 millones de filas y 1000 columnas en Databricks.
\end{itemize}

\subsection{Laboratorio: Ver Colab}
\begin{itemize}
    \item Scikit-learn: Regresión lineal, Regularización, Métricas
    \item Pipeline datos: EDA + Feature engineering + Regresión multiple + Regularizacion + metricas
\end{itemize}

\subsection{Anexo: Fundamento matematico y computacional}
\begin{itemize}
    \item Demostracion de la regresión lineal
    \item Demostracion de la regularización
    \item Optimizacion computacional
    \item Complejidad algoritmica    
\end{itemize}

% ------------------------------------------------------------------
% SEMANA 4: CLASIFICACIÓN I - REGRESIÓN LOGÍSTICA Y MÉTRICAS
% ------------------------------------------------------------------
\chapter{Semana 4: Clasificación I - Regresión Logística y Balance de Datos}

\section{Logro de la sesión}
Construir, interpretar y evaluar modelos de clasificación binaria y multiclase utilizando regresión logística, incorporando técnicas de balanceo de datos y selección adecuada de métricas.

\subsection{Problematica de negocio}
\begin{itemize}
    \item Tipos de problemas de clasificación: binaria y multiclase
    \item Impacto del desbalance de clases en problemas reales
    \item Ejemplos clasicos: detección de fraude, churn de clientes, diagnóstico médico, clasificación de productos y clientes, score crediticio
    \item Diferencias clave entre regresión y clasificación (output continuo vs categórico)
\end{itemize}

\subsection{Modelado de clasificación}
\begin{itemize}
    \item Requisitos del modelo
    \item Regresión logística binaria: función sigmoide
    \item Extensión a regresión logística multiclase (One-vs-Rest, Softmax)
    \item Interpretación probabilística de las predicciones
    \item Interpretación de coeficientes (log-odds) e impacto de variables
    \item Límite de decisión (threshold) y trade-off entre métricas
    \item Regularización: L1 (Lasso), L2 (Ridge) y Elastic Net
    \item Plantilla base de Python: Imports, Split de datos, Modelo, Entrenamiento, Predicciones, Evaluacion
\end{itemize}

\subsection{Balance de datos}
\begin{itemize}
    \item Problema de clases desbalanceadas
    \item Técnicas de balanceo:
    \begin{itemize}
        \item Oversampling (SMOTE)
        \item Undersampling
        \item Ajuste de pesos (class weights)
    \end{itemize}
    \item Cuándo usar cada técnica
    \item Impacto en el modelo y en las métricas
\end{itemize}

\subsection{Métricas}
\begin{itemize}
    \item Matriz de confusión e interpretacion: TP, FP, TN, FN 
    \item Metricas e interpretacion: Exactitud, Precisión, recall (sensibilidad), especificidad
    \item F1-score y su interpretación
    \item Curva ROC y AUC; PR-AUC
    \item Métricas en escenarios desbalanceados
    \item Selección de métricas según contexto de negocio (agua potable, phising, covid)
\end{itemize}

\subsection{Comunicacion}
\begin{itemize}
    \item Storytelling e insights con data dummy
    \item Interpretación de errores: costo de falsos positivos vs falsos negativos
    \item Ajuste del threshold según objetivos de negocio
    \item Elevator pitch a equipo tecnico y no tecnico
\end{itemize}

\subsection{Reto: 1 punto}
\begin{itemize}
    \item Analizar un dataset desbalanceado, comparar métricas antes y después de aplicar técnicas de balanceo, y justificar cuál estrategia es más adecuada según el caso de negocio.
\end{itemize}

\subsection{Laboratorio: Ver Colab}
\begin{itemize}
    \item Scikit-learn: regresión logística (binaria y multiclase)
    \item Pipeline completo:
    \begin{itemize}
        \item EDA
        \item Feature engineering
        \item Balanceo de datos
        \item Entrenamiento de modelo
        \item Evaluación con métricas
        \item Ajuste de threshold
    \end{itemize}
\end{itemize}

\subsection{Anexo: Fundamento matematico y computacional}
\begin{itemize}
    \item Derivación de la regresión logística
    \item Función de costo: log-loss (cross-entropy)
    \item Optimización: gradiente descendente
    \item Extensión a multiclase (softmax)
    \item Complejidad algorítmica
\end{itemize}

% ------------------------------------------------------------------
% SEMANA 5: CLASIFICACIÓN II - ALGORITMOS CLÁSICOS
% ------------------------------------------------------------------
\chapter{Semana 5: Clasificación II - Algoritmos Clásicos}

\section{Logro de la sesión}
Construir, comparar e interpretar modelos de clasificación utilizando algoritmos clásicos (KNN, Naive Bayes y SVM), comprendiendo sus supuestos, ventajas y limitaciones en distintos contextos de negocio.

\subsection{Problematica de negocio}
\begin{itemize}
    \item Selección de algoritmos de clasificación según el tipo de datos
    \item Comparación de desempeño entre modelos
    \item Casos donde la regresión logística no es suficiente (no linealidad, alta dimensionalidad)
    \item Impacto del tamaño de datos y escalamiento en el rendimiento
\item KNN: sistemas de recomendación simples, detección de similitud
    \item Naive Bayes: clasificación de texto (spam, sentimiento)
    \item SVM: problemas con fronteras no lineales (visión por computadora, bioinformática)
\end{itemize}

\subsection{Modelado}
\begin{itemize}
    \item K-Nearest Neighbors (KNN):
    \begin{itemize}
        \item Requisitos del modelo
        \item Concepto de distancia (euclidiana, manhattan)
        \item Elección de $k$ y trade-off bias-varianza
        \item Sensibilidad a la escala de variables
        \item Plantilla base de Python: Imports, Split de datos, Modelo, Entrenamiento, Predicciones, Evaluacion
    \end{itemize}
    
    \item Naive Bayes:
    \begin{itemize}
        \item Requisitos del modelo
        \item Teorema de Bayes
        \item Supuesto de independencia condicional
        \item Variantes: Gaussian, Multinomial, Bernoulli
        \item Interpretación probabilística
        \item Plantilla base de Python: Imports, Split de datos, Modelo, Entrenamiento, Predicciones, Evaluacion
    \end{itemize}
    
    \item Support Vector Machines (SVM):
    \begin{itemize}
        \item Requisitos del modelo
        \item Margen máximo y vectores de soporte
        \item Kernel lineal vs no lineal (RBF)
        \item Hiperparámetros: $C$ y $\gamma$
        \item Manejo de no linealidad
        \item Plantilla base de Python: Imports, Split de datos, Modelo, Entrenamiento, Predicciones, Evaluacion
    \end{itemize}
    
    \item Comparación conceptual entre algoritmos:
    \begin{itemize}
        \item Paramétricos vs no paramétricos
        \item Interpretabilidad vs performance
        \item Escalabilidad computacional
    \end{itemize}
\end{itemize}

\subsection{Métricas}
\begin{itemize}
    \item Matriz de confusión: TP, FP, TN, FN
    \item Exactitud, Precisión, recall, F1-score
    \item Curva ROC y AUC
    \item Comparación de modelos usando métricas consistentes
\end{itemize}

\subsection{Comunicacion}
\begin{itemize}
    \item Comparación de modelos para toma de decisiones
    \item Explicación de trade-offs: interpretabilidad vs precisión
    \item Justificación del modelo seleccionado según contexto de negocio
\end{itemize}

\subsection{Reto: 1 punto}
\begin{itemize}
    \item Comparar KNN, Naive Bayes y SVM en un mismo dataset, evaluando métricas y tiempos de entrenamiento, y justificar cuál modelo es más adecuado según el caso.
\end{itemize}

\subsection{Laboratorio}
\begin{itemize}
    \item Implementación de KNN, Naive Bayes y SVM en Scikit-learn
    \item Pipeline:
    \begin{itemize}
        \item Datos preprocesados (EDA + feature engineering ya realizado)
        \item Escalamiento de variables (clave para KNN y SVM)
        \item Entrenamiento de múltiples modelos
        \item Evaluación con métricas
        \item Comparación de resultados
    \end{itemize}
\end{itemize}

\subsection{Anexo: Fundamento matematico y computacional}
\begin{itemize}
    \item Distancias y espacios métricos (KNN)
    \item Teorema de Bayes y probabilidades condicionales
    \item Optimización en SVM (margen máximo)
    \item Complejidad computacional de cada algoritmo
\end{itemize}

% ------------------------------------------------------------------
% SEMANA 6: ÁRBOLES DE DECISIÓN Y RANDOM FOREST
% ------------------------------------------------------------------
\chapter{Semana 6: Árboles de Decisión y Random Forest}

\section{Logro de la sesión}
Construir, interpretar y comparar modelos basados en árboles de decisión y ensambles tipo Random Forest, comprendiendo su capacidad para capturar relaciones no lineales y su aplicación en problemas reales de clasificación y regresión.

\subsection{Problematica de negocio}
\begin{itemize}
    \item Modelos interpretables vs modelos de alto performance
    \item Manejo de relaciones no lineales y variables categóricas
    \item Reducción de overfitting en modelos complejos
    \item Necesidad de entender qué variables impactan más en la predicción
    \item Detección de fraude en tarjetas de crédito
    \item Scoring crediticio y evaluación de riesgo
    \item Modelos con alta dimensionalidad (texto, genómica)
    \item Problemas donde la interpretabilidad es clave (reglas de decisión)
\end{itemize}

\subsection{Modelado}
\begin{itemize}
    \item Árboles de decisión (clasificación y regresión):
    \begin{itemize}
        \item Estructura de árbol: nodos, ramas y hojas
        \item Criterios de división:
        \begin{itemize}
            \item Clasificación: Gini, entropía
            \item Regresión: MSE
        \end{itemize}
        \item Sobreajuste (overfitting) y control de complejidad:
        \begin{itemize}
            \item \texttt{max\_depth}, \texttt{min\_samples\_split}, \texttt{min\_samples\_leaf}
        \end{itemize}
        \item Interpretación de decisiones del modelo
        \item Plantilla base de Python: Imports, Split de datos, Modelo, Entrenamiento, Predicciones, Evaluacion
    \end{itemize}

    \item Random Forest (Bagging):
    \begin{itemize}
        \item Concepto de ensamble: combinación de múltiples árboles
        \item Bootstrap sampling (muestreo con reemplazo)
        \item Selección aleatoria de variables (feature randomness)
        \item Reducción de varianza vs árbol individual
        \item Out-of-Bag (OOB) error como validación interna
        \item Plantilla base de Python: Imports, Split de datos, Modelo, Entrenamiento, Predicciones, Evaluacion
    \end{itemize}

    \item Comparación:
    \begin{itemize}
        \item Árbol simple vs Random Forest
        \item Interpretabilidad vs performance
        \item Bias vs varianza
    \end{itemize}
\end{itemize}

\subsection{Métricas}
\begin{itemize}
    \item Clasificación: precisión, recall, F1-score, AUC-ROC
    \item Regresión: MSE, RMSE, MAE, $R^2$
    \item Importancia de variables:
    \begin{itemize}
        \item Gini importance (feature importance)
        \item Permutation importance
    \end{itemize}
    \item Selección de métricas según objetivo de negocio
\end{itemize}

\subsection{Comunicacion}
\begin{itemize}
    \item Explicación de decisiones del modelo mediante reglas del árbol
    \item Uso de importancia de variables para generar insights
    \item Comparación clara entre modelos simples y ensambles
    \item Traducción de resultados a impacto de negocio
\end{itemize}

\subsection{Reto: 1 punto}
\begin{itemize}
    \item Comparar un árbol de decisión y un Random Forest en el mismo dataset, analizando métricas, overfitting y la importancia de variables.
\end{itemize}

\subsection{Laboratorio}
\begin{itemize}
    \item Experimentacion: Implementación de árboles y Random Forest en Scikit-learn
    \item Pipeline:
    \begin{itemize}
        \item Datos preprocesados (EDA + feature engineering)
        \item Entrenamiento de árbol de decisión
        \item Control de overfitting (tuning de hiperparámetros)
        \item Entrenamiento de Random Forest
        \item Evaluación y comparación de modelos
        \item Análisis de importancia de variables
    \end{itemize}
\end{itemize}

\subsection{Anexo: Fundamento matematico y computacional}
\begin{itemize}
    \item Criterios de impureza: Gini, entropía, MSE
    \item Reducción de varianza en bagging
    \item Bias-variance trade-off
    \item Complejidad computacional de árboles y ensambles
\end{itemize}

% ------------------------------------------------------------------
% SEMANA 7: GRADIENT BOOSTING (XGBoost, LightGBM)
% ------------------------------------------------------------------
\chapter{Semana 7: Gradient Boosting (XGBoost, LightGBM)}

\section{Logro de la sesión}
Construir, optimizar e interpretar modelos de Gradient Boosting (XGBoost y LightGBM), comprendiendo su funcionamiento secuencial y ajustando hiperparámetros para maximizar el rendimiento en problemas de clasificación y regresión.

\subsection{Problematica de negocio}
\begin{itemize}
    \item Necesidad de modelos de alto rendimiento en datos tabulares
    \item Mejora incremental de predicciones (reducción de error)
    \item Manejo de relaciones no lineales y variables complejas
    \item Trade-off entre performance, tiempo de entrenamiento e interpretabilidad
    \item Predicción de clics (CTR) en publicidad online
    \item Scoring crediticio y detección de fraude
    \item Modelado de datos tabulares estructurados
\end{itemize}

\subsection{Modelado}
\begin{itemize}
    \item Boosting:
    \begin{itemize}
        \item Ensamble secuencial de modelos débiles
        \item Corrección iterativa de errores
        \item Diferencia clave con bagging (Random Forest)
    \end{itemize}

    \item Gradient Boosting:
    \begin{itemize}
        \item Optimización de una función de pérdida (loss function)
        \item Uso de árboles débiles (shallow trees)
        \item Learning rate (shrinkage) y número de estimadores
        \item Riesgo de overfitting
    \end{itemize}

    \item XGBoost:
    \begin{itemize}
        \item Requisitos del modelo
        \item Regularización (L1, L2)
        \item Manejo de valores nulos
        \item Early stopping
        \item Paralelización eficiente
        \item Plantilla base de Python: Imports, Split de datos, Modelo, Entrenamiento, Predicciones, Evaluacion
    \end{itemize}

    \item LightGBM:
    \begin{itemize}
        \item Requisitos del modelo
        \item Gradient-based One-Side Sampling (GOSS)
        \item Histogram-based splitting
        \item Manejo eficiente de variables categóricas
        \item Entrenamiento más rápido en grandes volúmenes
        \item Plantilla base de Python: Imports, Split de datos, Modelo, Entrenamiento, Predicciones, Evaluacion
    \end{itemize}

    \item Comparación:
    \begin{itemize}
        \item Random Forest vs Gradient Boosting (bagging vs boosting)
        \item XGBoost vs LightGBM
        \item Interpretabilidad vs performance
    \end{itemize}
\end{itemize}

\subsection{Métricas}
\begin{itemize}
    \item Clasificación: log-loss, AUC-ROC, F1-score
    \item Regresión: RMSE, MAE, $R^2$
    \item Evaluación en validación cruzada
    \item Selección de métricas según el objetivo del negocio
\end{itemize}

\subsection{Hiperparámetros clave}
\begin{itemize}
    \item \texttt{n\_estimators}: número de árboles
    \item \texttt{learning\_rate}: tamaño del paso en cada iteración
    \item \texttt{max\_depth}: complejidad de los árboles
    \item \texttt{subsample}, \texttt{colsample\_bytree}: control de overfitting
    \item Regularización: \texttt{lambda}, \texttt{alpha} (XGBoost)
\end{itemize}

\subsection{Comunicacion}
\begin{itemize}
    \item Explicación de mejoras de performance frente a modelos base
    \item Interpretación de importancia de variables
    \item Justificación del uso de modelos complejos
    \item Balance entre precisión y explicabilidad
\end{itemize}

\subsection{Reto: 1 punto}
\begin{itemize}
    \item Comparar Random Forest vs XGBoost vs LightGBM en un dataset de kaggle, optimizando hiperparámetros y evaluando diferencias en performance y tiempo de entrenamiento.
\end{itemize}

\subsection{Laboratorio}
\begin{itemize}
    \item Implementación de XGBoost y LightGBM
    \item Pipeline:
    \begin{itemize}
        \item Datos preprocesados (EDA + feature engineering)
        \item Entrenamiento de modelo base
        \item Ajuste de hiperparámetros
        \item Uso de early stopping
        \item Evaluación y comparación de modelos
        \item Análisis de importancia de variables
    \end{itemize}
\end{itemize}

\subsection{Anexo: Fundamento matematico y computacional}
\begin{itemize}
    \item Gradient Boosting como optimización de funciones
    \item Expansión funcional (additive models)
    \item Regularización en boosting
    \item Complejidad computacional y escalabilidad
\end{itemize}

% ------------------------------------------------------------------
% SEMANA 8: EVALUACIÓN Y VALIDACIÓN DE MODELOS
% ------------------------------------------------------------------
\chapter{Semana 8: Evaluación y Validación de Modelos}

\section{Logro de la sesión}
Evaluar, validar y seleccionar modelos de machine learning de forma robusta, asegurando su capacidad de generalización mediante técnicas de validación cruzada y optimización de hiperparámetros.

\subsection{Problematica de negocio}
\begin{itemize}
    \item Modelos que funcionan bien en entrenamiento pero fallan en producción
    \item Selección del mejor modelo entre múltiples alternativas
    \item Riesgo de sobreajuste en datasets pequeños o complejos
    \item Necesidad de estimar performance real antes de desplegar
\end{itemize}

\subsection{Aplicaciones y casos típicos}
\begin{itemize}
    \item Selección de modelos en proyectos reales
    \item Validación antes de despliegue en producción
\end{itemize}

\subsection{Modelado y validación}
\begin{itemize}
    \item Underfitting y overfitting:
    \begin{itemize}
        \item Identificación y diagnóstico
        \item Estrategias de mitigación
    \end{itemize}

    \item Bias-Variance tradeoff:
    \begin{itemize}
        \item Interpretación práctica
        \item Relación con complejidad del modelo
    \end{itemize}

    \item Validación de modelos:
    \begin{itemize}
        \item Train / Validation / Test split
        \item Validación cruzada:
        \begin{itemize}
            \item K-Fold
            \item Stratified K-Fold (clasificación)
            \item Leave-One-Out
        \end{itemize}
    \end{itemize}

    \item Curvas de aprendizaje:
    \begin{itemize}
        \item Diagnóstico de underfitting vs overfitting
        \item Interpretación de curvas
    \end{itemize}
\end{itemize}

\subsection{Optimización de modelos}
\begin{itemize}
    \item Búsqueda de hiperparámetros:
    \begin{itemize}
        \item Grid Search
        \item Randomized Search
        \item Introducción a Bayesian Optimization
    \end{itemize}
    
    \item Buenas prácticas:
    \begin{itemize}
        \item Evitar data leakage
        \item Uso de pipelines
        \item Separación correcta de datos
    \end{itemize}
\end{itemize}

\subsection{Métricas}
\begin{itemize}
    \item Uso de métricas en validación cruzada
    \item Clasificación: AUC, F1-score, log-loss
    \item Regresión: RMSE, MAE, $R^2$
    \item Selección de métricas alineadas al objetivo de negocio
    \item Comparación de modelos basada en métricas promedio
\end{itemize}

\subsection{Plantilla base de Python}
\begin{itemize}
    \item \texttt{sklearn.model\_selection}:
    \begin{itemize}
        \item \texttt{train\_test\_split}
        \item \texttt{cross\_val\_score}
        \item \texttt{StratifiedKFold}
        \item \texttt{learning\_curve}
        \item \texttt{GridSearchCV}
        \item \texttt{RandomizedSearchCV}
    \end{itemize}
\end{itemize}

\subsection{Comunicacion}
\begin{itemize}
    \item Explicación de performance esperada en producción
    \item Comparación clara entre modelos (baseline vs optimizado)
    \item Justificación de selección de modelo
    \item Comunicación de riesgos (overfitting, variabilidad)
\end{itemize}

\subsection{Reto: 1 punto}
\begin{itemize}
    \item Implementar validación cruzada y búsqueda de hiperparámetros en un modelo (por ejemplo, Random Forest o XGBoost), comparando resultados antes y después del tuning.
\end{itemize}

\subsection{Laboratorio}
\begin{itemize}
    \item Pipeline completo de validación:
    \begin{itemize}
        \item Train/test split
        \item Validación cruzada
        \item Entrenamiento de modelo base
        \item Búsqueda de hiperparámetros
        \item Evaluación final en test
    \end{itemize}
    \item Comparación de múltiples modelos (logística, árboles, boosting)
\end{itemize}

\subsection{Anexo: Fundamento matematico y computacional}
\begin{itemize}
    \item Bias-variance decomposition
    \item Fundamento de validación cruzada
    \item Complejidad computacional del tuning
    \item Riesgo de overfitting en selección de modelos
\end{itemize}

% ------------------------------------------------------------------
% SEMANA 9: APRENDIZAJE NO SUPERVISADO - CLUSTERING Y PCA
% ------------------------------------------------------------------
\chapter{Semana 9: Aprendizaje No Supervisado - Clustering y PCA}

\section{Logro de la sesión}
Explorar y extraer valor de datos no etiquetados mediante técnicas de clustering y reducción de dimensionalidad (PCA), identificando patrones, segmentos y estructuras ocultas en los datos.

\subsection{Problematica de negocio}
\begin{itemize}
    \item Falta de etiquetas en los datos (no hay variable objetivo)
    \item Necesidad de segmentar clientes o comportamientos
    \item Reducción de complejidad en datasets de alta dimensionalidad
    \item Identificación de patrones ocultos y anomalías
\end{itemize}

\subsection{Aplicaciones y casos típicos}
\begin{itemize}
    \item Segmentación de clientes (marketing)
    \item Reducción de dimensionalidad para visualización
    \item Compresión de datos
    \item Detección de anomalías (fraude, comportamiento atípico)
\end{itemize}

\subsection{Modelado}
\begin{itemize}
    \item PCA (Análisis de Componentes Principales):
    \begin{itemize}
        \item Reducción de dimensionalidad basada en varianza
        \item Componentes principales y su interpretación
        \item Explained variance ratio
        \item Visualización (2D/3D) para exploración
    \end{itemize}

    \item K-Means:
    \begin{itemize}
        \item Requisitos del modelo
        \item Algoritmo basado en centroides
        \item Inicialización y convergencia
        \item Elección de $k$:
        \begin{itemize}
            \item Método del codo
            \item Coeficiente de silueta
        \end{itemize}
        \item Sensibilidad a escala y outliers
    \end{itemize}

    \item DBSCAN:
    \begin{itemize}
        \item Requisitos del modelo
        \item Clustering basado en densidad
        \item Parámetros: $\epsilon$ (eps) y \texttt{min\_samples}
        \item Detección de ruido (outliers)
        \item Ventajas frente a K-Means (no requiere $k$)
    \end{itemize}

    \item Clustering jerárquico:
    \begin{itemize}
        \item Requisitos del modelo
        \item Enfoque aglomerativo
        \item Dendrogramas e interpretación
        \item Selección del número de clusters
    \end{itemize}

    \item Comparación de métodos:
    \begin{itemize}
        \item K-Means vs DBSCAN vs Jerárquico
        \item Escenarios de uso según tipo de datos
    \end{itemize}
\end{itemize}

\subsection{Métricas}
\begin{itemize}
    \item Coeficiente de silueta
    \item Índice de Davies-Bouldin
    \item Inercia (K-Means)
    \item Limitaciones de métricas sin ground truth
\end{itemize}

\subsection{Python}
\begin{itemize}
    \item \texttt{sklearn.decomposition.PCA}
    \item \texttt{sklearn.cluster.KMeans}, \texttt{DBSCAN}, \texttt{AgglomerativeClustering}
    \item \texttt{sklearn.metrics.silhouette\_score}
    \item Buenas prácticas:
    \begin{itemize}
        \item Escalamiento de datos antes de clustering
        \item Uso de PCA para visualización
    \end{itemize}
\end{itemize}

\subsection{Comunicacion}
\begin{itemize}
    \item Interpretación de clusters (perfilamiento)
    \item Traducción de segmentos a decisiones de negocio
    \item Visualización de resultados (scatter plots, PCA)
    \item Limitaciones del análisis no supervisado
\end{itemize}

\subsection{Reto: 1 punto}
\begin{itemize}
    \item Aplicar K-Means y DBSCAN a un dataset real, comparar resultados y construir perfiles de los clusters identificados.
\end{itemize}

\subsection{Laboratorio}
\begin{itemize}
    \item Pipeline de análisis no supervisado:
    \begin{itemize}
        \item Datos preprocesados (EDA + feature engineering)
        \item Escalamiento de variables
        \item Reducción de dimensionalidad (PCA)
        \item Clustering (K-Means, DBSCAN)
        \item Evaluación con métricas
        \item Visualización e interpretación de clusters
    \end{itemize}
\end{itemize}

\subsection{Anexo: Fundamento matematico y computacional}
\begin{itemize}
    \item Álgebra lineal en PCA (autovalores y autovectores)
    \item Optimización en K-Means
    \item Concepto de densidad en DBSCAN
    \item Complejidad computacional de algoritmos de clustering
\end{itemize}

% ------------------------------------------------------------------
% SEMANA 10: SERIES DE TIEMPO
% ------------------------------------------------------------------
\chapter{Semana 10: Series de Tiempo}

\section{Logro de la sesión}
Modelar y pronosticar datos temporales utilizando técnicas estadísticas clásicas, machine learning y librerías modernas, extrayendo patrones de tendencia, estacionalidad y ciclos.

\subsection{Problematica de negocio}
\begin{itemize}
    \item Pronóstico de demanda y ventas para planificación
    \item Detección de patrones estacionales y tendencias
    \item Análisis de comportamiento temporal para optimización de recursos
    \item Integración de modelos clásicos y machine learning para predicción
\end{itemize}

\subsection{Aplicaciones y casos típicos}
\begin{itemize}
    \item Pronóstico de ventas en retail
    \item Predicción de demanda energética o consumo de servicios
    \item Análisis y predicción de tráfico web o logs de sistemas
    \item Forecast financiero (ingresos, precios)
\end{itemize}

\subsection{Modelado}
\begin{itemize}
    \item Componentes de series temporales:
    \begin{itemize}
        \item Tendencia, estacionalidad, ciclos y ruido
        \item Descomposición aditiva y multiplicativa
    \end{itemize}
    
    \item Feature engineering temporal:
    \begin{itemize}
        \item Rezagos (lags)
        \item Ventanas móviles (rolling statistics)
        \item Diferencias (differencing)
        \item Variables de calendario (día de semana, mes, festivos)
    \end{itemize}
    
    \item Modelos clásicos:
    \begin{itemize}
        \item ARIMA (AutoRegressive Integrated Moving Average)
        \item SARIMA (Seasonal ARIMA)
        \item Requisitos del modelo
        \item Evaluación de residuales, estacionariedad y autocorrelación
    \end{itemize}
    
    \item Modelos modernos:
    \begin{itemize}
        \item Prophet (Facebook) para tendencias y estacionalidad flexibles
        \item Requisitos del modelo
        \item Modelos basados en machine learning:
        \begin{itemize}
            \item Regresión con características temporales (lags, rolling features)
            \item XGBoost o LightGBM con características de tiempo
        \end{itemize}
    \end{itemize}
    
    \item Comparación:
    \begin{itemize}
        \item Modelos estadísticos vs modelos ML
        \item Precisión vs interpretabilidad
        \item Escenarios de uso según tipo de serie y granularidad
    \end{itemize}
\end{itemize}

\subsection{Métricas}
\begin{itemize}
    \item MAE (Mean Absolute Error)
    \item MAPE (Mean Absolute Percentage Error)
    \item RMSE (Root Mean Squared Error)
    \item MASE (Mean Absolute Scaled Error)
    \item Selección de métricas según negocio (por ejemplo, errores relativos vs absolutos)
\end{itemize}

\subsection{Python}
\begin{itemize}
    \item \texttt{pandas} para manejo de fechas, creación de lags y rolling windows
    \item \texttt{statsmodels} para ARIMA, SARIMA y descomposición de series
    \item \texttt{fbprophet} (Prophet) para modelado flexible de series
    \item Machine learning: \texttt{sklearn}, \texttt{xgboost}, \texttt{lightgbm} aplicados a series con features temporales
    \item Evaluación de modelos con métricas: MAE, MAPE, RMSE, MASE
\end{itemize}

\subsection{Comunicacion}
\begin{itemize}
    \item Visualización de tendencias y estacionalidades
    \item Comparación de pronósticos vs datos reales
    \item Explicación del impacto de patrones estacionales y ciclos en decisiones de negocio
    \item Comunicación de incertidumbre y límites del modelo
\end{itemize}

\subsection{Reto: 1 punto}
\begin{itemize}
    \item Construir un pronóstico de ventas con ARIMA y Prophet, comparar métricas de error y explicar la elección de features temporales.
\end{itemize}

\subsection{Laboratorio}
\begin{itemize}
    \item Pipeline de series temporales:
    \begin{itemize}
        \item Preparación de datos y creación de lags
        \item Descomposición de series para análisis de tendencia y estacionalidad
        \item Entrenamiento de modelos ARIMA/SARIMA
        \item Entrenamiento de Prophet y/o modelo ML con features temporales
        \item Evaluación de predicciones con métricas
        \item Visualización de pronósticos vs valores reales
    \end{itemize}
\end{itemize}

\subsection{Anexo: Fundamento matematico y computacional}
\begin{itemize}
    \item Autocorrelación y estacionariedad
    \item Descomposición aditiva y multiplicativa
    \item Optimización de parámetros ARIMA/SARIMA (p,d,q)
    \item Concepto de errores y métricas escaladas
    \item Complejidad computacional en series temporales y modelos ML
\end{itemize}

% ------------------------------------------------------------------
% SEMANA 11: MODELOS COMPLEMENTARIOS (REGRESIÓN, CLASIFICACIÓN, CLUSTERING)
% ------------------------------------------------------------------
\chapter{Semana 11: Modelos Complementarios}

\section{Logro de la sesión}
Conocer y aplicar algoritmos adicionales de regresión y clustering que son útiles para problemas específicos o entrevistas técnicas, comprendiendo cuándo y por qué utilizarlos.

\subsection{Problematica de negocio}
\begin{itemize}
    \item Manejo de datos con outliers o distribuciones no normales
    \item Modelado de relaciones no lineales (regresión polinómica)
    \item Segmentación de clientes o elementos con estructuras complejas o jerárquicas
    \item Detección de patrones arbitrarios en datasets no etiquetados
\end{itemize}

\subsection{Aplicaciones y casos típicos}
\begin{itemize}
    \item Regresión robusta: predicción de precios de vivienda con outliers
    \item SVR: problemas de regresión no lineales con pocos datos
    \item DBSCAN: detección de anomalías, clusters irregulares
    \item Clustering jerárquico: segmentación de clientes, análisis filogenético
\end{itemize}

\subsection{Modelado}
\begin{itemize}
    \item Regresión polinómica:
    \begin{itemize}
        \item Extensión de la regresión lineal
        \item Introducción de términos polinómicos
        \item Riesgo de overfitting y regularización
    \end{itemize}

    \item Regresión robusta:
    \begin{itemize}
        \item HuberRegressor: combina MSE y MAE para outliers
        \item RANSACRegressor: ajuste iterativo ignorando outliers
        \item Aplicaciones en precios extremos o datos ruidosos
    \end{itemize}

    \item SVM para regresión (SVR):
    \begin{itemize}
        \item Función de pérdida epsilon-insensitive
        \item Kernel lineal y no lineal
        \item Regularización y parámetro C
    \end{itemize}

    \item Clustering avanzado:
    \begin{itemize}
        \item DBSCAN: clustering basado en densidad, detección de ruido
        \item Clustering jerárquico: dendrogramas, aglomerativo vs divisivo
        \item Elección de parámetros y análisis de estructura jerárquica
    \end{itemize}

    \item Métricas adicionales:
    \begin{itemize}
        \item Clustering: Adjusted Rand Index, información mutua
        \item Interpretación de silueta y Davies-Bouldin para clusters complejos
    \end{itemize}
\end{itemize}



\subsection{Métricas}
\begin{itemize}
    \item Regresión robusta: MAE, RMSE, $R^2$
    \item Clustering: Silhouette, Davies-Bouldin, Adjusted Rand Index, Mutual Information
\end{itemize}

\subsection{Python}
\begin{itemize}
    \item Regresión robusta: \texttt{sklearn.linear\_model.HuberRegressor}, \texttt{RANSACRegressor}
    \item SVR: \texttt{sklearn.svm.SVR}
    \item Clustering avanzado: \texttt{sklearn.cluster.DBSCAN}, \texttt{AgglomerativeClustering}
    \item Métricas: \texttt{sklearn.metrics.adjusted\_rand\_score}, \texttt{silhouette\_score}, \texttt{davies\_bouldin\_score}
\end{itemize}

\subsection{Comunicacion}
\begin{itemize}
    \item Interpretación de modelos robustos frente a outliers
    \item Visualización y explicación de clusters complejos
    \item Comunicación clara de limitaciones y supuestos de cada técnica
\end{itemize}

\subsection{Reto: 1 punto}
\begin{itemize}
    \item Aplicar RANSAC y HuberRegressor en un dataset con outliers y comparar resultados con regresión lineal estándar.
    \item Aplicar DBSCAN y clustering jerárquico a un dataset sin etiquetas y comparar la estructura de clusters.
\end{itemize}

\subsection{Laboratorio}
\begin{itemize}
    \item Pipeline de modelos complementarios:
    \begin{itemize}
        \item Selección de dataset con outliers o estructura compleja
        \item Entrenamiento de regresión robusta (Huber y RANSAC)
        \item Entrenamiento de SVR
        \item Aplicación de DBSCAN y clustering jerárquico
        \item Evaluación de métricas y visualización de resultados
    \end{itemize}
\end{itemize}

\subsection{Anexo: Fundamento matematico y computacional}
\begin{itemize}
    \item Funciones de pérdida robusta y su derivación
    \item Optimización iterativa en RANSAC
    \item Kernels y SVR
    \item Complejidad de clustering densidad vs jerárquico
\end{itemize}

% ------------------------------------------------------------------
% SEMANA 12: SISTEMAS DE RECOMENDACIÓN
% ------------------------------------------------------------------
\chapter{Semana 12: Sistemas de Recomendación}

\section{Logro de la sesión}
Diseñar e implementar sistemas de recomendación basados en filtrado colaborativo y factorización matricial, evaluando su desempeño y entendiendo limitaciones como cold start y sparsity.

\subsection{Problematica de negocio}
\begin{itemize}
    \item Personalización de productos y contenidos para usuarios
    \item Manejo de datos dispersos (sparse data) y gran escala
    \item Problema de cold start: nuevos usuarios o nuevos ítems
    \item Selección de métricas adecuadas para distintos objetivos (clasificación vs predicción de rating)
\end{itemize}

\subsection{Aplicaciones y casos típicos}
\begin{itemize}
    \item Recomendación de productos en Amazon
    \item Películas en Netflix
    \item Música en Spotify
    \item Contenidos de noticias o artículos personalizados
\end{itemize}

\subsection{Modelado}
\begin{itemize}
    \item Filtrado colaborativo:
    \begin{itemize}
        \item Basado en usuarios (user-user)
        \item Basado en ítems (item-item)
        \item Medidas de similitud: coseno, correlación de Pearson
    \end{itemize}
    
    \item Factorización matricial:
    \begin{itemize}
        \item SVD (Singular Value Decomposition)
        \item Funk SVD (gradient descent sobre ratings)
        \item Reducción de dimensionalidad para embeddings de usuarios e ítems
        \item Interpretación de factores latentes
    \end{itemize}
    
    \item Problemas comunes:
    \begin{itemize}
        \item Cold start: nuevos usuarios o ítems sin historial
        \item Sparsity: matriz de ratings muy dispersa
        \item Escalabilidad para grandes datasets
    \end{itemize}
    
    \item Métricas de evaluación:
    \begin{itemize}
        \item Ranking: Precision@K, Recall@K, NDCG
        \item Predicción de ratings: RMSE, MAE
        \item Tradeoff entre precisión y diversidad de recomendaciones
    \end{itemize}
\end{itemize}

\subsection{Métricas}
\begin{itemize}
    \item Ranking: Precision@K, Recall@K, NDCG
    \item Predicción de rating: RMSE, MAE
    \item Evaluación comparativa de modelos de filtrado colaborativo vs factorización matricial
\end{itemize}

\subsection{Python}
\begin{itemize}
    \item Librería \texttt{surprise}:
    \begin{itemize}
        \item \texttt{SVD}, \texttt{KNNBasic}, \texttt{KNNWithMeans}
        \item Split train/test y cross-validation
        \item Cálculo de RMSE, MAE, Precision@K
    \end{itemize}
    \item Alternativa con \texttt{sklearn.decomposition.TruncatedSVD} para factorización matricial
    \item Embeddings con redes neuronales (opcional, avanzado)
    \item Escalamiento de matrices dispersas (\texttt{scipy.sparse}) para datasets grandes
\end{itemize}

\subsection{Comunicacion}
\begin{itemize}
    \item Interpretación de factores latentes y similitudes entre usuarios/ítems
    \item Explicación de recomendaciones a stakeholders
    \item Visualización de top-k recomendaciones por usuario o ítem
    \item Limitaciones del modelo: cold start y sparsity
\end{itemize}

\subsection{Reto: 1 punto}
\begin{itemize}
    \item Implementar un sistema de recomendación con SVD sobre un dataset de ratings, evaluar Precision@K y RMSE, y proponer estrategias para cold start.
\end{itemize}

\subsection{Laboratorio}
\begin{itemize}
    \item Pipeline de recomendación:
    \begin{itemize}
        \item Preparación de datos de usuarios e ítems
        \item Filtrado colaborativo (user-user, item-item)
        \item Factorización matricial con SVD
        \item Evaluación con métricas de ranking y predicción
        \item Visualización de recomendaciones personalizadas
    \end{itemize}
\end{itemize}

\subsection{Anexo: Fundamento matematico y computacional}
\begin{itemize}
    \item Descomposición en valores singulares y gradiente descendente en Funk SVD
    \item Cálculo de similitud entre vectores de usuarios/ítems
    \item Complejidad computacional para matrices dispersas
    \item Estrategias de optimización y escalabilidad
\end{itemize}
% ------------------------------------------------------------------
% SEMANA 13: INTERPRETABILIDAD DE MODELOS
% ------------------------------------------------------------------
\chapter{Semana 13: Interpretabilidad de Modelos}

\section{Logro de la sesión}
Explicar las predicciones de modelos complejos mediante técnicas de interpretabilidad global y local, permitiendo comprender decisiones de negocio, identificar sesgos y cumplir regulaciones.

\subsection{Problematica de negocio}
\begin{itemize}
    \item Modelos complejos (árboles, ensembles, boosting, deep learning) difíciles de interpretar.
    \item Necesidad de justificar decisiones ante stakeholders o regulaciones.
    \item Identificación de características influyentes y posibles sesgos.
\end{itemize}

\subsection{Aplicaciones y casos típicos}
\begin{itemize}
    \item Justificación de decisiones a stakeholders (ej. rechazo de crédito).
    \item Cumplimiento de regulaciones como GDPR (“derecho a explicación”).
    \item Depuración de modelos: detectar sesgos y variables irrelevantes.
    \item Comunicación de insights a equipos de negocio y técnicos.
\end{itemize}

\subsection{Modelado}
\begin{itemize}
    \item Importancia de características:
    \begin{itemize}
        \item Modelos basados en árboles: Gini importance, permutación.
        \item Evaluación de impacto global de cada variable.
    \end{itemize}
    
    \item Partial Dependence Plots (PDP) e Individual Conditional Expectation (ICE):
    \begin{itemize}
        \item PDP: efecto promedio de una característica sobre la predicción.
        \item ICE: efecto individual de la característica para cada observación.
    \end{itemize}
    
    \item SHAP (SHapley Additive exPlanations):
    \begin{itemize}
        \item Concepto de valores de Shapley para cada feature.
        \item Interpretación local (predicción individual) y global (importancia promedio).
        \item Compatible con modelos complejos como XGBoost, LightGBM o redes neuronales.
    \end{itemize}
    
    \item LIME (Local Interpretable Model-agnostic Explanations):
    \begin{itemize}
        \item Explicación local aproximando el modelo por uno interpretable.
        \item Útil para explicar predicciones individuales sin conocer internals del modelo.
    \end{itemize}
    
    \item Diferencias clave:
    \begin{itemize}
        \item Global vs local: PDP/feature importance vs SHAP/LIME
        \item Model-agnostic vs específico del modelo
    \end{itemize}
\end{itemize}

\subsection{Métricas}
\begin{itemize}
    \item No existen métricas estándar; la evaluación se centra en consistencia, coherencia y utilidad de las explicaciones.
    \item Comparación de interpretaciones globales vs locales para validar insights.
\end{itemize}

\subsection{Python}
\begin{itemize}
    \item \texttt{sklearn.inspection}: \texttt{plot\_partial\_dependence}, \texttt{PartialDependenceDisplay} para PDP/ICE.
    \item \texttt{shap}: cálculo de valores SHAP, visualización de summary plots, dependence plots, force plots.
    \item \texttt{lime}: \texttt{LimeTabularExplainer} para explicaciones locales.
    \item Buenas prácticas: usar PDP y SHAP conjuntamente para ver coherencia global y local.
\end{itemize}

\subsection{Comunicacion}
\begin{itemize}
    \item Visualización clara de la importancia de features y su impacto.
    \item Explicación de predicciones individuales mediante SHAP o LIME.
    \item Identificación de posibles sesgos y recomendaciones de mitigación.
    \item Traducción de insights técnicos a lenguaje de negocio.
\end{itemize}

\subsection{Reto: 1 punto}
\begin{itemize}
    \item Tomar un modelo de clasificación o regresión entrenado previamente, generar explicaciones con SHAP y LIME, y analizar la consistencia de las conclusiones.
\end{itemize}

\subsection{Laboratorio}
\begin{itemize}
    \item Pipeline de interpretabilidad:
    \begin{itemize}
        \item Cargar modelo entrenado (árboles, ensemble o boosting)
        \item Calcular importancia de features global
        \item Generar PDP e ICE para características clave
        \item Aplicar SHAP y LIME para explicaciones locales
        \item Visualizar resultados y generar reporte interpretativo para stakeholders
    \end{itemize}
\end{itemize}

\subsection{Anexo: Fundamento matematico y computacional}
\begin{itemize}
    \item Teoría de valores de Shapley y propiedades (eficiencia, simetría, aditividad)
    \item Algoritmos de aproximación de SHAP y complejidad
    \item Optimización de cálculos PDP y ICE en grandes datasets
    \item Consideraciones computacionales para LIME (subsampling, kernel)
\end{itemize}

% ------------------------------------------------------------------
% SEMANA 14: DESPLIEGUE DE MODELOS (MLOps)
% ------------------------------------------------------------------
\chapter{Semana 14: Despliegue de Modelos (MLOps)}

\section{Logro de la sesión}
Llevar un modelo a producción, comprender los conceptos clave de MLOps, monitorear su rendimiento y aplicar pruebas A/B para validar mejoras.

\subsection{Problematica de negocio}
\begin{itemize}
    \item Separación entre entrenamiento y producción para evitar errores.
    \item Necesidad de versionar modelos y datos para reproducibilidad.
    \item Detectar degradación de modelos debido a cambios en los datos.
    \item Evaluar mejoras de nuevos modelos mediante experimentos controlados (A/B testing).
\end{itemize}

\subsection{Aplicaciones y casos típicos}
\begin{itemize}
    \item Integrar un sistema de recomendación en una app web
    \item Actualizar modelos de scoring crediticio periódicamente
    \item Detectar degradación de modelos predictivos por cambios en comportamiento del usuario
    \item Test A/B para evaluar nuevos algoritmos de recomendación o predicción
\end{itemize}

\subsection{Modelado y despliegue}
\begin{itemize}
    \item Separación entrenamiento vs inferencia:
    \begin{itemize}
        \item Entrenamiento offline, inferencia online
        \item Pipeline reproducible con preprocesamiento y transformación de features
    \end{itemize}

    \item Serialización de modelos:
    \begin{itemize}
        \item Formatos: \texttt{pickle}, \texttt{joblib}, \texttt{ONNX} para interoperabilidad
    \end{itemize}
    
    \item Creación de APIs:
    \begin{itemize}
        \item Flask o FastAPI para endpoints de predicción
        \item Validación de entrada y manejo de errores
    \end{itemize}
    
    \item Contenerización:
    \begin{itemize}
        \item Docker para portar modelos y dependencias
        \item Reproducibilidad en distintos entornos
    \end{itemize}

    \item Versionamiento:
    \begin{itemize}
        \item Modelos: MLflow, DVC
        \item Datos y features: Git + DVC
        \item Registro de experimentos y métricas
    \end{itemize}
    
    \item Monitoreo:
    \begin{itemize}
        \item Performance en producción: RMSE, AUC, precisión, recall según el modelo
        \item Data drift y concept drift: tests estadísticos (Kolmogorov-Smirnov, Population Stability Index)
        \item Alertas automáticas si métricas caen por debajo de umbrales
    \end{itemize}
    
    \item Pruebas A/B:
    \begin{itemize}
        \item Comparación de modelos en producción con tráfico real
        \item Medición de métricas clave (conversión, click-through, engagement)
        \item Decisión basada en resultados estadísticamente significativos
    \end{itemize}
\end{itemize}

\subsection{Métricas}
\begin{itemize}
    \item Métricas tradicionales según tipo de modelo: RMSE, MAE, AUC, precisión, recall
    \item Drift de datos y predicciones: Kolmogorov-Smirnov, PSI
    \item Resultados de experimentos A/B: lift, significancia estadística, tasa de conversión
\end{itemize}

\subsection{Python}
\begin{itemize}
    \item Serialización: \texttt{pickle}, \texttt{joblib}, exportación a \texttt{ONNX}
    \item Creación de APIs: Flask o FastAPI con carga de modelo y endpoints de predicción
    \item Versionamiento y tracking: \texttt{MLflow}, \texttt{DVC}
    \item Monitoreo de drift: \texttt{evidently}, \texttt{alibi-detect}
    \item Implementación básica de test A/B: dividir tráfico, recolectar métricas y comparar resultados estadísticos
\end{itemize}

\subsection{Comunicacion}
\begin{itemize}
    \item Explicar diferencias de performance entre modelos a stakeholders
    \item Reportar resultados de tests A/B con métricas claras y gráficas
    \item Alertas y dashboards para seguimiento de modelos en producción
\end{itemize}

\subsection{Reto: 1 punto}
\begin{itemize}
    \item Desplegar un modelo de clasificación o regresión en Flask/FastAPI, realizar monitoreo de métricas y ejecutar un test A/B para comparar con la versión anterior.
\end{itemize}

\subsection{Laboratorio}
\begin{itemize}
    \item Pipeline de despliegue:
    \begin{itemize}
        \item Entrenamiento y serialización del modelo
        \item Creación de API con Flask/FastAPI
        \item Contenerización con Docker
        \item Versionamiento con MLflow/DVC
        \item Monitoreo de métricas y detección de drift
        \item Ejecución de test A/B sobre un grupo de usuarios
    \end{itemize}
\end{itemize}

\subsection{Anexo: Fundamento matematico y computacional}
\begin{itemize}
    \item Serialización y compatibilidad entre entornos
    \item Versionamiento de modelos y control de datasets
    \item Test estadístico A/B: significancia, intervalos de confianza
    \item Optimización de pipelines de inferencia para baja latencia
\end{itemize}

% ------------------------------------------------------------------
% METODOLOGÍA Y EVALUACIÓN
% ------------------------------------------------------------------
\chapter{Metodología y Evaluación}

\section{Estrategias metodológicas}
\begin{itemize}
    \item \textbf{Clases síncronas:} Exposición conceptual con diapositivas, resolución de problemas en pizarra y programación en vivo con Python utilizando cuadernos de Jupyter.
    \item \textbf{Trabajo asíncrono:} Lectura de material complementario, resolución de ejercicios prácticos y pequeños proyectos semanales.
    \item \textbf{Aprendizaje activo:} Discusión de casos reales, ejercicios de entrevistas técnicas y retos de código.
    \item \textbf{Recursos:} Plataforma virtual, Zoom, grabaciones, foros y cuadernos de Jupyter.
\end{itemize}

\section{Materiales educativos}
\begin{itemize}
    \item Presentaciones interactivas (PPT, videos).
    \item Cuadernos de Jupyter con ejemplos y ejercicios.
    \item Bibliografía: libros clásicos (Géron, Hastie) y documentación oficial.
    \item Google Colab para ejecución sin instalación.
\end{itemize}

\section{Evaluación}
\begin{table}[h]
\centering
\begin{tabular}{lc}
\toprule
\textbf{Ítem} & \textbf{Ponderación} \\
\midrule
Evaluaciones continuas (EC) & 40\% \\
Examen parcial (EP) & 30\% \\
Examen final (EF) & 30\% \\
\midrule
\textbf{Promedio final (PF)} & PF = 0.4 \times EC + 0.3 \times EP + 0.3 \times EF \\
\bottomrule
\end{tabular}
\caption{Sistema de evaluación}
\end{table}
\begin{itemize}
    \item \textbf{Evaluaciones continuas:} Controles de lectura, entrega de ejercicios de programación, participación.
    \item \textbf{Examen parcial:} Cubre semanas 1-7.
    \item \textbf{Examen final:} Cubre semanas 8-14.
    \item \textbf{Examen sustitutorio:} Semana 15, reemplaza la nota más baja entre parcial y final.
\end{itemize}

% ------------------------------------------------------------------
% BIBLIOGRAFÍA
% ------------------------------------------------------------------
\chapter*{Bibliografía}
\addcontentsline{toc}{chapter}{Bibliografía}

\begin{itemize}
    \item Géron, A. (2019). \textit{Hands-On Machine Learning with Scikit-Learn, Keras, and TensorFlow}. O'Reilly.
    \item Hastie, T., Tibshirani, R., Friedman, J. (2009). \textit{The Elements of Statistical Learning}. Springer.
    \item Bishop, C. M. (2006). \textit{Pattern Recognition and Machine Learning}. Springer.
    \item VanderPlas, J. (2016). \textit{Python Data Science Handbook}. O'Reilly.
    \item Documentación oficial de scikit-learn, XGBoost, LightGBM, SHAP, etc.
\end{itemize}

\vspace{1cm}
\begin{flushright}
\textit{Elaborado por:} Carlos César Sánchez Coronel \\
\textit{Aprobado por:} Dirección de Ingeniería de Inteligencia Artificial \\
\textit{Fecha:} 17/01/2026
\end{flushright}

\end{document}